\chapter{Σύνολο Δεδομένων}

Το παρόν κεφάλαιο τεκμηριώνει το σύνολο δεδομένων MAVOS-DD που χρησιμοποιείται
στην εργασία, τη διαδικασία επιλογής υποσυνόλου, τη λογική διαχωρισμού train/val/test
και τους ελέγχους ακεραιότητας που εφαρμόστηκαν σε κάθε φάση.

% ─────────────────────────────────────────────────────────────────────────────
\section{Το Dataset MAVOS-DD}
% ─────────────────────────────────────────────────────────────────────────────

Το MAVOS-DD (\textit{Multimodal Audio-Visual Open-Set Deepfake Detection}) είναι
ένα δημόσιο σύνολο δεδομένων για την αξιολόγηση συστημάτων ανίχνευσης deepfake
σε βίντεο πολλαπλών πηγών. Περιλαμβάνει αυθεντικά βίντεο από YouTube καθώς
και συνθετικά βίντεο που έχουν παραχθεί από τέσσερις γεννήτριες κινούμενου
προσώπου, οργανωμένα σε υποφακέλους ανά πηγή παραγωγής (\codeid{source\_folder}).
Είναι διαθέσιμο μέσω του Hugging Face Hub ως \codeid{unibuc-cs/MAVOS-DD}.

\textbf{Γιατί επιλέχθηκε.} Το MAVOS-DD παρέχει πολλαπλές \emph{οπτικά παραποιημένες}
πηγές (γεννήτριες κινούμενου προσώπου) με \emph{αυθεντικό ηχητικό κανάλι},
κάτι που επιτρέπει τον ξεχωριστό έλεγχο του visual branch χωρίς σύγχυση από
ηχητική παραποίηση. Ταυτόχρονα, η διαθεσιμότητά του μέσω Hub επιτρέπει
streaming λήψη με ελεγχόμενα per-class ανώτατα όρια, αποφεύγοντας
την αποθήκευση ολόκληρου του dataset.

\textbf{Αγγλόφωνο υποσύνολο.} Επιλέγουμε αποκλειστικά βίντεο με
\codeid{language = english}, ώστε τα TTS spoof δεδομένα (ElevenLabs, Google Neural2,
OpenAI) να αντιστοιχούν στη γλώσσα των αυθεντικών μεταγραφών. Εγγραφές
σε άλλες γλώσσες απορρίπτονται σε επίπεδο μεταδεδομένων κατά την ενσωμάτωση.

\textbf{Φάκελοι πηγής (source folders).} Το dataset οργανώνεται
στις παρακάτω κατηγορίες:

\begin{itemize}
  \item \codeid{real} — Αυθεντικά βίντεο με πραγματικό πρόσωπο και πραγματικό ήχο
    από YouTube. Δυαδική ετικέτα: 0 (bonafide).
  \item \codeid{echomimic} — Συνθετικά βίντεο από την γεννήτρια EchoMimic:
    κινούμενο πρόσωπο από εικόνα, με πραγματικό ήχο οδηγό.
    Δυαδική ετικέτα: 1 (fake).
  \item \codeid{memo} — Συνθετικά βίντεο από τη γεννήτρια MEMO:
    ανάλογη αρχή με EchoMimic, διαφορετική αρχιτεκτονική παραγωγής.
    Δυαδική ετικέτα: 1 (fake).
  \item \codeid{liveportrait} — Συνθετικά βίντεο από τη γεννήτρια LivePortrait.
    Εισάγεται στο Phase~6+. Δυαδική ετικέτα: 1 (fake).
  \item \codeid{sonic} — Συνθετικά βίντεο από τη γεννήτρια Sonic.
    Εισάγεται στο Phase~6+. Δυαδική ετικέτα: 1 (fake).
\end{itemize}

Σε κάθε περίπτωση, ο \codeid{binary\_label} ορίζεται ως
$0$ για \codeid{real} και $1$ για όλες τις υπόλοιπες πηγές,
σύμφωνα με τον χάρτη ετικετών \codeid{LABEL\_MAP} στο \codeid{src/common.py}.
Σημείωση σχεδιασμού: ο ήχος στα \codeid{echomimic}/\codeid{memo}/\codeid{liveportrait}/\codeid{sonic}
βίντεο είναι \emph{αυθεντικός} — η παραποίηση αφορά αποκλειστικά το οπτικό κανάλι.

% ─────────────────────────────────────────────────────────────────────────────
\section{Αρχική Επιλογή — Phase~1–5 (1\,000 βίντεο)}
% ─────────────────────────────────────────────────────────────────────────────

Για τις Phases 1–5 εφαρμόστηκε ένα ανώτατο όριο 1\,000 βίντεο συνολικά,
κατανεμημένων σε τρεις πηγές:

\begin{center}
\begin{tabular}{lr}
  \toprule
  \codeid{source\_folder} & Ανώτατο όριο (cap) \\
  \midrule
  \codeid{real}      & 500 \\
  \codeid{echomimic} & 250 \\
  \codeid{memo}      & 250 \\
  \midrule
  Σύνολο             & 1\,000 \\
  \bottomrule
\end{tabular}
\end{center}

\textbf{Σκεπτικό επιλογής ορίου.} Το cap των 1\,000 βίντεο επιλέχθηκε
για τις Phases 1–5 ως smoke-test / πρώιμο baseline:
αρκετό δείγμα για να επιβεβαιωθεί η ορθότητα του pipeline
(ingestion, εξαγωγή χαρακτηριστικών, εκπαίδευση, αξιολόγηση)
χωρίς την πλήρη φόρτωση του dataset, η οποία απαιτεί σημαντικούς
πόρους αποθήκευσης και χρόνο εξαγωγής χαρακτηριστικών.
Οι πηγές \codeid{liveportrait} και \codeid{sonic} δεν συμπεριλήφθηκαν
στο αρχικό cap επειδή η ανάγκη για οπτικά διαφοροποιημένες πηγές
ανακαλύφθηκε ως ανάγκη αφότου τα Phase~5 αποτελέσματα έδειξαν
δομική κατάρρευση του visual branch (βλ. Κεφάλαιο~6).

% ─────────────────────────────────────────────────────────────────────────────
\section{Επέκταση Phase~6+ — \textasciitilde{}4\,149 βίντεο}
% ─────────────────────────────────────────────────────────────────────────────

Από το Phase~6 και εξής, το ανώτατο όριο επεκτάθηκε σε περίπου 4\,149 βίντεο
σε πέντε πηγές, σύμφωνα με την Revision~1 του roadmap
(\codeid{docs/roadmap-audio-visual-speech-detection.md}).
Ο πίνακας~\ref{tab:caps} συγκρίνει τα ανώτατα όρια των δύο φάσεων.

\begin{table}[ht]
  \centering
  \begin{tabular}{lrr}
    \toprule
    \texttt{source\_folder} & Phase 1–5 & Phase 6+ \\
    \midrule
    real         & 500 & 2{,}500 \\
    echomimic    & 250 &   600 \\
    memo         & 250 &   400 \\
    liveportrait &   - &   314 \\
    sonic        &   - &   335 \\
    \midrule
    Σύνολο       & 1{,}000 & 4{,}149 \\
    \bottomrule
  \end{tabular}
  \caption{Όρια ανά \texttt{source\_folder} για τις δύο φάσεις.}
  \label{tab:caps}
\end{table}

\textbf{Αιτιολόγηση επέκτασης.} Δύο λόγοι κινητοποίησαν την αύξηση:

\begin{enumerate}
  \item \textbf{Δομική αποτυχία visual branch:} Στο Phase~5 κάθε bonafide
    εγγραφή και η αντίστοιχη spoof εγγραφή της μοιράζονταν
    τον ίδιο \codeid{source\_video\_id} και, κατ'~επέκταση,
    τα ίδια \codeid{.npz} χαρακτηριστικά χειλιών.
    Ένα BiGRU που καλείται να διακρίνει δύο σειρές με
    byte-πανομοιότυπη είσοδο δεν μπορεί να κάνει καλύτερα
    από τη συχνότητα κλάσης (val ROC-AUC = 0.57,
    confusion matrix $\mathrm{tn}=0 / \mathrm{fp}=150 / \mathrm{fn}=0 / \mathrm{tp}=217$).
    Η προσθήκη \codeid{liveportrait} και \codeid{sonic} εισάγει
    σειρές όπου τα ίδια τα lip features διαφέρουν μεταξύ
    πραγματικού και ψεύτικου προσώπου.
  \item \textbf{Μεγαλύτερο bonafide pool:} Η αύξηση από 500 σε 2\,500
    πραγματικά βίντεο ελέγχει αν οι ανιχνευτές εκμεταλλεύονται
    συντόμευσεις συνθηκών εγγραφής αντί για αληθινά
    διακριτικά χαρακτηριστικά.
\end{enumerate}

\textbf{Ιδιότητα ιδεμποτεντότητας.} Το \codeid{src/data/download\_subset.py}
εκτελείται ασφαλώς πολλαπλές φορές: επαναφορτώνει τις
υπάρχουσες μετρήσεις ανά κλάση από το \codeid{data/manifest.csv}
και κατεβάζει μόνο τα βίντεο που λείπουν έως τα νέα caps.
Δεν πραγματοποιείται διπλή λήψη.

% ─────────────────────────────────────────────────────────────────────────────
\section{Stratified Frozen Split}
% ─────────────────────────────────────────────────────────────────────────────

Οι διαχωρισμοί train/val/test υπολογίστηκαν μία φορά και παγοποιήθηκαν
με το αρχείο \codeid{src/data/make\_splits.py}. Η στρατηγική:

\begin{itemize}
  \item \textbf{Αναλογία:} 70\,\% εκπαίδευση / 15\,\% επικύρωση / 15\,\% δοκιμή.
  \item \textbf{Στρωτοποίηση (stratification):} πάνω στη στήλη \codeid{source\_folder},
    ώστε κάθε split να διατηρεί την κατανομή των πηγών.
  \item \textbf{Seed:} 42 παντού, για πλήρη αναπαραγωγιμότητα.
  \item \textbf{Έξοδος:} \codeid{data/splits/train.csv}, \codeid{data/splits/val.csv},
    \codeid{data/splits/test.csv}.
\end{itemize}

\textbf{Το script είναι destructive.} Το \codeid{make\_splits.py} γράφει
άνευ όρων στα \codeid{data/splits/*.csv} χωρίς έλεγχο ύπαρξης αρχείου.
Πριν την επέκταση Phase~6+ (re-freeze στα νέα caps), έγινε εφεδρικό
αντίγραφο (backup) των Phase~1–5 splits. Τα checkpoints των Phases~1–5
παραμένουν δεσμευμένα στα αρχικά splits και αξιολογούνται αποκλειστικά
με αυτά.

\textbf{Κανόνας αξιολόγησης test split.} Το test split αξιολογείται
\emph{ακριβώς μία φορά} ανά checkpoint, μετά την επιλογή μοντέλου
βάσει val metrics. Η εντολή \codeid{python -m src.evaluate --checkpoint
<ckpt> --split test --allow-test} απαιτεί ρητή σημαία
\codeid{--allow-test} ακριβώς για να αποτρέπει τυχαία
πρόωρη αξιολόγηση.

% ─────────────────────────────────────────────────────────────────────────────
\section{Επιμέλεια και Έλεγχοι Ακεραιότητας}
% ─────────────────────────────────────────────────────────────────────────────

Μετά την ενσωμάτωση, εκτελείται ο ελεγκτής ακεραιότητας μέσω της εντολής:

\begin{center}
  \codeid{python -m src.data.download\_subset --validate}
\end{center}

Η λειτουργία αυτή είναι read-only: δεν κατεβάζει ούτε τροποποιεί δεδομένα.
Ελέγχει τα παρακάτω:

\begin{itemize}
  \item \textbf{Πλήθος εγγραφών:} Ο αριθμός γραμμών στο \codeid{data/manifest.csv}
    πρέπει να ισούται με το άθροισμα των caps.
  \item \textbf{Ανά κλάση caps:} Κάθε \codeid{source\_folder} πρέπει να έχει
    ακριβώς \codeid{CAPS[source\_folder]} εγγραφές.
  \item \textbf{Μοναδικότητα video\_id:} Δεν επιτρέπονται διπλότυπα
    \codeid{video\_id} στο manifest.
  \item \textbf{Παρουσία αρχείων:} Κάθε \codeid{relative\_path} πρέπει να αντιστοιχεί
    σε υπαρκτό αρχείο \codeid{.mp4} στο \codeid{data/raw/}.
  \item \textbf{Συνέπεια ετικετών:} Ο \codeid{binary\_label} κάθε εγγραφής
    πρέπει να συμφωνεί με το \codeid{LABEL\_MAP[source\_folder]}.
  \item \textbf{Μη μηδενικά πεδία video probe:} Τα πεδία \codeid{duration\_s},
    \codeid{fps} και \codeid{n\_frames} πρέπει να είναι θετικοί αριθμοί.
\end{itemize}

\textbf{Quarantine.} Κατά την ενσωμάτωση, κάθε βίντεο που αποτυγχάνει
στους ελέγχους (αδύνατη ανάγνωση, μηδενικά frames, μηδενικό fps,
απουσία ηχητικής ροής) μετακινείται στο \codeid{data/quarantine/<source\_folder>/}
και καταγράφεται στο \codeid{data/quarantine\_log.csv} με τον λόγο απόρριψης.
Κανένα βίντεο δεν αποκλείεται αθόρυβα.
